\runninghead{The Sunday and Its Texts}
\section{The Resumed Sixth Sunday after the Epiphany and Its Texts}
\label{sec:opening}

The Gospel of this Mass gives two short parables that Jesus tells the crowds
from a boat on the Sea of Galilee (Mt 13:1--2). The first is of a mustard seed
that a man sows in his field, ``the least indeed of all seeds; but when it is
grown up, it is greater than all herbs, and becometh a tree, so that the birds
of the air come, and dwell in the branches thereof.'' The second is of leaven,
``which a woman took and hid in three measures of meal, until the whole was
leavened.'' Matthew then remarks that Jesus spoke to the crowds only in
parables, and quotes a psalm as fulfilled in that way of speaking: ``I will
open my mouth in parables, I will utter things hidden from the foundation of
the world'' (Mt 13:31--35). Neither parable is explained in the chapter. When
the disciples come to Jesus in the house afterwards they ask him only about
the cockle sown among the wheat (Mt 13:36), and St John Chrysostom supposes
that they let the mustard seed and the leaven pass ``as being plainer''
(\work{Hom.\ in Mt.}~47.1).

The Epistle is the opening of St Paul's first letter to the Thessalonians.
Paul had preached in their synagogue for three sabbaths; some believed, a mob
attacked the house of Jason where he lodged, and the brethren sent him away by
night to Berea (Acts 17:1--10). He writes to the young church remembering ``the
work of your faith and labour and charity: and \dots\ the enduring of the
hope.'' The gospel came to them ``not \dots\ in word only, but in power also,''
and they received the word ``in much tribulation, with joy of the Holy Ghost.''
From them ``was spread abroad the word of the Lord not only in Macedonia and in
Achaia but also in every place.'' Everyone, he says, reports how they ``turned
to God from idols to serve the living and true God. And to wait for his Son
from heaven \dots\ Jesus, who hath delivered us from the wrath to come'' (1~Thess
1:2--10).

The five chants come from the Twenty-third Sunday after Pentecost. The Introit
gives God's words through Jeremias to the exiles in Babylon: ``I think thoughts
of peace, and not of affliction; you shall call upon Me and I will hear you;
and I will bring back your captivity from all places,'' with the psalm verse
``Lord, Thou hast blessed Thy land: Thou hast turned away the captivity of
Jacob'' (Jer 29:11--12, 14; Ps 84:2). The Gradual thanks God for a deliverance:
``Thou hast delivered us, O Lord, from them that afflict us: and hast put them
to shame that hate us. In God we will glory all the day'' (Ps 43:8--9). The
Alleluia and then the Offertory sing ``From the depths I have cried to Thee, O
Lord: Lord, hear my prayer'' (Ps 129:1--2), and the Offertory repeats its first
line. The Communion is the Lord's promise in Mark: ``whatsoever you ask when you
pray, believe that you shall receive, and it shall be done to you'' (Mk 11:24).
The English of the chants is that of the Lasance hand missal; the lessons are
given in the Douay--Rheims.

The three prayers of the priest are short. The Collect asks that, ``ever
fixing our thoughts on reasonable things, we may both in word and in deed do
what is pleasing to Thee.'' The Secret asks that this offering ``cleanse,
renew, govern, and protect us.'' The Postcommunion, said by those who have been
``fed with celestial delights,'' asks that ``we may ever hunger after those
things by which we truly live.''

The Missal prints this Mass in the Proper of Time after the Twenty-third
Sunday after Pentecost, under the heading \latin{Dominica sexta quae superfuit
post Epiphaniam}, the sixth Sunday ``left over'' after the Epiphany, of the
second class (printed pp.~421--422, nos.~1759--1768). Its Collect, Epistle,
Gospel, Secret and Postcommunion are those of the Sixth Sunday after the
Epiphany; its Introit, Gradual, Alleluia, Offertory and Communion are those of
the Twenty-third Sunday after Pentecost, as RGMR~298 directs for every Sunday
after the Epiphany resumed between the twenty-third and the last Sunday after
Pentecost. In 2026 Septuagesima fell on 1~February and left room for only three
Sundays after the Epiphany. There are twenty-six Sundays after Pentecost this
year, and in that case RG~18~b sets the Fifth Sunday after the Epiphany in the
twenty-fourth place and the Sixth in the twenty-fifth, which falls on
15~November; the last Sunday after Pentecost follows on 22~November. The
calendar's feast of St Albert the Great on 15~November, of the third class, is
not commemorated, because a Sunday of the second class admits only the
commemoration of a feast of the second class (RG~111~b). The Mass is said in
green, with the Gloria and the Credo, and with the Preface of the Most Holy
Trinity, which the Missal directs after the Secret (RGMR~494~b).

Each of the three interpretations that follow takes up one question. The
Gospel raises the first two. The first is how a seed that is the least of all
becomes a tree. St Hilary of Poitiers says that the seed is the Lord himself, whose
strength ``is kindled by tribulations and pressures'' and who was sown when he
was seized, put to death and buried. St Ambrose, expounding the same parable in
Luke, says that mustard is a cheap and simple thing that gives out its strength
when it is ground, and that faith does the same when adversity presses it. St
John Chrysostom says the parable was given to disciples who feared that they
were too few. The first interpretation holds that the church at Thessalonica,
which received the word in much tribulation and whose faith then went out into
every place, was pressed and spread in the way Ambrose describes.

The second question is what the leaven is, what the three measures are, and
how far the leavening goes. St Augustine says the leaven is love, ``which
warms and rouses,'' and the measures heart, soul and mind; St Bede, in Luke,
that love hidden in the mind must grow until it has changed the whole mind;
St Jerome, that the leaven is the knowledge of the Scriptures, working until
spirit, soul and body agree; St Ambrose, that it is the Church's spiritual
teaching, working in flesh, soul and spirit until the whole inner man is
leavened. The second interpretation takes the charity that Paul remembers in
the Thessalonians, beside the work of their faith and their labour, for the
love that Augustine and Bede call the leaven, at work in a church; and it
reads the Collect's petition for word and deed by the leaven's work in the
whole person.

The third question comes from the tenses of the Mass. The Introit's verse says
``Thou hast turned away the captivity,'' the Gradual ``Thou hast delivered
us,'' and the Epistle names Jesus ``who hath delivered us from the wrath to
come''; yet the antiphon promises a return, Paul's converts wait for the Son
from heaven. St Augustine says that the prophets speak of the future as past because
it is as certain as if it were past, and that the psalmist prays for what he
has just called done ``because they are not yet done.'' St Jerome says that
God could have kept his promise to the exiles without their prayer, and urged
them to pray so that they might deserve to receive it. The third interpretation reads
the Epistle's waiting for the Son and the Communion's promise by those
answers: what God has promised is certain, and his people pray and wait for
it.

\begin{maptable}
Introit &
\latin{Dicit Dominus: Ego cogito cogitationes pacis}; Jer 29:11--12, 14;
verse Ps 84:2 (85:2) &
The Lord thinks thoughts of peace and not of affliction; you shall call upon
him and he will hear; he will bring back your captivity from all places; the
psalm says he has blessed his land and turned away the captivity of Jacob.\\
Collect &
\latin{Praesta, quaesumus, omnipotens Deus: ut, semper rationabilia
meditantes} &
That, ever fixing our thoughts on reasonable things, we may do what pleases
God in word and in deed.\\
Epistle &
1~Thess 1:2--10 &
Paul thanks God for the Thessalonians' work of faith, labour, charity and
enduring hope; the gospel came to them in power; they received the word in
much tribulation with joy of the Holy Ghost, and their faith is spoken of
everywhere; they turned from idols to serve the living God and to wait for his
Son, who has delivered us from the wrath to come.\\
Gradual &
\latin{Liberasti nos, Domine, ex affligentibus nos}; Ps 43:8--9 (44:8--9) &
God has delivered us from those who afflict us and put to shame those who
hate us; in God we will glory all the day and praise his name for ever.\\
Alleluia &
\latin{De profundis clamavi ad te, Domine}; Ps 129:1--2 (130:1--2) &
From the depths I have cried to the Lord; Lord, hear my prayer.\\
Gospel &
Mt 13:31--35 &
The kingdom is like a mustard seed, least of all seeds, that becomes a tree
for the birds, and like leaven hidden in three measures of meal until the
whole is leavened; Jesus speaks to the crowds only in parables, fulfilling the
psalm's ``I will utter things hidden from the foundation of the world.''\\
Offertory &
\latin{De profundis clamavi ad te, Domine}; Ps 129:1--2 (130:1--2) &
The cry from the depths sung again, its first line repeated.\\
Secret &
\latin{Haec nos oblatio, Deus, mundet} &
That this offering cleanse, renew, govern and protect us.\\
Communion &
\latin{Amen dico vobis, quidquid orantes petitis}; Mk 11:24 &
Whatever you ask when you pray, believe that you shall receive, and it shall
be done for you.\\
Postcommunion &
\latin{Caelestibus, Domine, pasti deliciis} &
That, fed with heavenly delights, we may always long for the things by which
we truly live.\\
\end{maptable}
